\chapter{Anaesthetics: An Overview}
\index{Anaesthetics: An Overview}

\section*{Definition and Overview}
Anaesthetics are the medicines and techniques used to make surgery and other procedures pain-free and safe. Anaesthesia can be general, in which the person is placed into a reversible state of unconsciousness with no awareness or memory, or regional and local, in which part of the body is numbed while the person remains awake or lightly sedated. The specialty also covers pain relief during and after procedures and around childbirth.

Modern anaesthesia is delivered by specialist doctors and anaesthetic nurses who continuously monitor breathing, heart, blood pressure, oxygen levels, and depth of consciousness throughout the procedure. Anaesthesia is a partnership: the anaesthetist plans a tailored anaesthetic for each person based on the procedure, the person's health, and their preferences, and adjusts it moment by moment in the operating theatre.

\section*{Epidemiology}
Anaesthesia accompanies a very large share of hospital care. Millions of general anaesthetics are given each year in developed countries alone, and anaesthesia-related care is delivered alongside most surgery, many dental procedures, and childbirth. Serious complications directly attributable to anaesthesia are now rare, with major adverse events occurring in the order of a few per thousand or fewer.

The risk varies with the urgency and complexity of the surgery, the person's age and general health, and the type of anaesthetic used. Older people, those with heart or lung disease, and those having emergency surgery carry higher risk, all of which is assessed and planned for before the operation.

\section*{Aetiology and Pathophysiology}
Anaesthetics work by reversibly suppressing nervous-system activity. Their mechanisms are complex and not fully understood, but the main principles are:
\begin{itemize}
    \item \textbf{General anaesthesia:} inhaled gases or intravenous drugs produce loss of consciousness, pain relief, muscle relaxation, and amnesia, acting mainly on the brain and spinal cord
    \item \textbf{Regional anaesthesia:} injections near nerve bundles (spinal and epidural) or around large nerves (nerve blocks) interrupt pain signals from a region of the body
    \item \textbf{Local anaesthesia:} creams or injections numb a small area by blocking sodium channels in nerve endings
    \item \textbf{Sedation:} medicines that relax and reduce awareness while preserving breathing, used alone or with local or regional techniques
\end{itemize}
The anaesthetist selects the safest combination for each person, taking into account the type of surgery, the person's medical conditions, and the medications they take.

\section*{Clinical Features}
Before an anaesthetic, the assessment covers everything the anaesthetist needs to plan care; after a procedure, normal recovery has recognisable features:
\begin{itemize}
    \item \textbf{Pre-operative:} discussion of past and current health, medicines, allergies, and any previous anaesthetic problems
    \item \textbf{During the procedure:} continuous monitoring of breathing, oxygen levels, heart rhythm, blood pressure, and depth of anaesthesia
    \item \textbf{After general anaesthesia:} drowsiness, nausea, a sore throat, dizziness, and feeling cold or briefly confused on waking
    \item \textbf{After regional anaesthesia:} numbness and weakness of the affected limb or region for several hours, with gradual return of normal feeling
    \item \textbf{Postoperative pain:} controlled with a combination of medicines and, where used, regional blocks
\end{itemize}

\section*{Differential Diagnosis}
Complications and normal recovery after anaesthesia can be confused with other problems:
\begin{itemize}
    \item Postoperative nausea versus nausea from opioid pain relief or from the surgery itself
    \item A sore throat from the breathing tube versus a developing postoperative chest infection
    \item Prolonged drowsiness versus low blood sugar or a neurological event
    \item Pain at the injection site versus bleeding or infection
    \item Anxiety before surgery versus a genuine medical complication
\end{itemize}

\section*{When to Seek Medical Care}
Before surgery, tell the anaesthetist about any chest pain, palpitations, breathlessness, recent illness, or concerns about pregnancy. After an anaesthetic, seek medical advice if:
\begin{itemize}
    \item Nausea, vomiting, or pain is severe or not settling
    \item You develop a fever, shaking, or a cough after surgery
    \item A regional block has not worn off within the expected time
    \item You have new numbness, tingling, or weakness in a limb
    \item You feel confused, drowsy, or are having difficulty breathing
\end{itemize}

\textbf{Contact your local emergency services for:}
\begin{itemize}
    \item Severe chest pain, sudden breathlessness, or loss of consciousness
    \item Sudden facial or limb weakness, difficulty speaking, or confusion (possible stroke)
\end{itemize}

\section*{Diagnosis and Investigations}
\begin{itemize}
    \item \textbf{Pre-anaesthetic assessment:} a detailed history, examination, and discussion of options, benefits, and risks
    \item \textbf{Blood tests:} full blood count, kidney and liver function, and blood sugar, as indicated
    \item \textbf{ECG and heart tests:} for people with heart disease or risk factors
    \item \textbf{Chest X-ray and breathing tests:} for people with lung disease or smoking history
    \item \textbf{Airway assessment:} examination of the mouth, teeth, and neck to plan airway management
    \item \textbf{Fasting and medication review:} clear instructions about when to stop eating and drinking and which medicines to take on the day
\end{itemize}

\section*{Management}
\begin{itemize}
    \item \textbf{General anaesthesia:} intravenous induction, maintenance with inhaled or intravenous agents, and a controlled reversal and wake-up at the end
    \item \textbf{Regional techniques:} spinal, epidural, or nerve blocks, sometimes combined with sedation
    \item \textbf{Local anaesthesia with or without sedation:} for minor procedures
    \item \textbf{Airway management:} a breathing tube or airway device, with mechanical ventilation as needed
    \item \textbf{Monitoring and safety:} continuous observation of vital signs, with standard safety checklists and trained personnel throughout
    \item \textbf{Pain relief after surgery:} multimodal analgesia combining medicines with regional blocks where appropriate
    \item \textbf{Recovery care:} observation in a recovery area until the person is awake, comfortable, and stable
\end{itemize}

\section*{Complications}
\begin{itemize}
    \item Nausea, vomiting, sore throat, and dizziness after waking
    \item Pain or bruising at injection and drip sites
    \item Headache after spinal or epidural anaesthesia
    \item Low blood pressure during or after the anaesthetic
    \item Awareness during anaesthesia, which is very rare with modern monitoring
    \item Allergic reactions to anaesthetic drugs, also rare
    \item Damage to teeth, lips, or the voice box during airway management
    \item Serious rare events including heart attack, stroke, aspiration of stomach contents, and malignant hyperthermia
    \item Prolonged numbness, weakness, or nerve irritation after a regional block
\end{itemize}

\section*{Prognosis and Outcomes}
Modern anaesthesia is very safe, and serious complications are uncommon in healthy people having planned surgery. Recovery from the effects of the anaesthetic is usually rapid, with most people alert and comfortable within hours, although some feel drowsy or nauseated for up to a day. The main risks to overall recovery come from the surgery itself and the person's underlying health rather than from the anaesthetic. People with heart, lung, or kidney disease and those having complex or emergency surgery carry higher risk, which the anaesthetist assesses and plans for in advance.

\section*{Prevention and Self-Care}
\begin{itemize}
    \item Be honest and complete at the pre-anaesthetic assessment, including medicines, herbal products, and recreational drugs
    \item Follow fasting instructions exactly to reduce the risk of aspiration
    \item Keep taking prescribed medicines as directed, unless told otherwise
    \item Stop smoking well before surgery, as it reduces breathing complications
    \item Ask questions until you understand your options and what to expect
    \item After the anaesthetic, arrange for someone to take you home if advised
    \item Do not drive, operate machinery, or make important decisions for 24 hours after a general anaesthetic
    \item Tell the anaesthetist about any previous anaesthetic problems, family reactions, or allergies
\end{itemize}

\section*{Sources and Further Reading}
The following resources provide reliable, peer-reviewed and patient-orientated information on anaesthesia and anaesthetic care.

\begin{itemize}
    \item NHS. \textit{Overview of anaesthesia, including general and regional anaesthetic.} https://www.nhs.uk/conditions/
    \item Mayo Clinic. \textit{Guide to anaesthesia types and what to expect.} https://www.mayoclinic.org/
    \item MSD Manual. \textit{Professional and patient information on anaesthesia and its risks.} https://www.msdmanuals.com/
    \item MedlinePlus. \textit{Health information on anaesthesia and anaesthetic medicines.} https://medlineplus.gov/
    \item NICE. \textit{Clinical guidance on perioperative care and patient safety.} https://www.nice.org.uk/
\end{itemize}
